\begin{table*}[t]
    \centering
    \resizebox{\textwidth}{!}{
    \setlength{\tabcolsep}{4pt}
    \begin{tabular}{l cc cccccccc}
    \toprule
     & \multicolumn{2}{c}{Evolutionary VCV} & \multicolumn{8}{c}{Intraspecific VCV (full model)} \\
    \cmidrule(lr){2-3} \cmidrule(lr){4-11}
    \rotatebox{45}{Analysis}
        & \rotatebox{45}{Full model}
        & \rotatebox{45}{Taxon means model}
        & \rotatebox{45}{Modern human}
        & \rotatebox{45}{Neandertal}
        & \rotatebox{45}{\textit{Pan paniscus}}
        & \rotatebox{45}{\textit{Pan troglodytes}}
        & \rotatebox{45}{\textit{G. beringei}}
        & \rotatebox{45}{\textit{G. gorilla}}
        & \rotatebox{45}{\textit{Po. abelii}}
        & \rotatebox{45}{\textit{Po. pygmaeus}} \\
    \midrule
    C\textsubscript{1} & 8.68 & 12.03 & 78.84 & 44.66 & 90.35 & 58.68 & 38.31 & 26.86 & 26.24 & 45.33 \\
    \midrule
    I\textsuperscript{2} & 16.76 & 23.3 & 75.82 & 35.85 & 22.47 & 62.97 & - & - & - & - \\
    \bottomrule
    \end{tabular}}
    \caption{Fréchet variance for each VCV posterior distribution. The evolutionary VCV is compared
    between the full model, which accounts for intraspecific variation, and the taxon means model, in
    which species means are fixed at their observed sample means and only the evolutionary VCV is
    estimated; the remaining columns give the intraspecific VCV inferred for each taxon under the full
    model, which has no analog under the taxon means model. Higher values indicate greater uncertainty
    in the inferred VCV matrix.}
    \label{table:frechetVariances}
\end{table*}
